\section{Introduction} 
%Describe OSPF
Interconnection networks consist of groups of nodes that carry aggregated traffic across multiple sites and capacity links. The task of network routing is to transmit data between sources and destinations. Within the same domain, this type of network is called an intra-domain network. Intra-domain networks require routing protocols to facilitate packet transmission between nodes. A widely used protocol in these networks is the Open Shortest Path First (OSPF) routing protocol \citep{rfc2328}. OSPF is a link-state protocol that detects topological changes within the Autonomous System (AS), routes packets based on the IP header, and determines the optimal path using Dijkstra’s algorithm.

%Describe OSPF Limitation
OSPF has limited flexibility in intra-domain routing. The authoritative OSPF standard uses a single configuration cost that restricts protocol adjustments to the link weights \citep{rfc2328}. OSPF does not account for the real-time condition of links between nodes. The protocol ignores real-time load and capacity when using static weights. These limitations can create network congestion \citep{1039866}. As a result, intra-domain routing faces challenges in maintaining connectivity because standard OSPF offers limited traffic-engineering capability.

%Describe Research Gap  
As network traffic demand grows, there is a need for better network routing solutions. Multi-Agent Reinforcement Learning with a Graph Neural Network (GNN) method has recently been applied to dynamic routing in mobile ad-hoc networks  \citep{11134361}. However, existing research mainly focuses on wireless, mobile, and topology-varying environments. Their applicability to intra-domain routing, where routing is driven by capacity and traffic, remains unexplored. Furthermore, other previous studies also evaluated routing performance using synthetic traffic rather than real network demand.

%Describe Proposed Approach
To address these gaps, this work applies and adapts Multi-Agent Reinforcement Learning with a Graph Neural Network (GNN) backbone to intra-domain telecommunication backbone routing. Reinforcement Learning is a machine learning paradigm in which agents learn optimal actions through a system of rewards and punishment \citep{SHAKARAMI2020107496}. Multi-agent Reinforcement Learning involves a set of agents interacting with each other in the environment, each doing trial and error to maximize reward \citep{marl-book}. Combining with the GNNs, which are designed to process and analyze graph structure data \citep{10960451}, the network topology can be effectively represented as a graph.

%Describe Contribution
This study evaluates the performance of Multi-Agent Reinforcement Learning (MARL) with Graph Neural Network (GNN) integration against Open Shortest Path First (OSPF). The evaluation is conducted on three distinct telecommunication backbone topologies to enable a comprehensive comparison. We use the Survivable Network Design Library \citep{SNDlib10, OrlowskiPioroTomaszewskiWessaely2010} as the resource for topology and traffic data. Regarding MARL training and execution, this work adopts a centralized training and decentralized execution framework. The work employs analytical surrogate methods for training due to resource constraints and performs packet-level evaluations using NS-3 \citep{Riley2010}. Furthermore, this work measures performance by collecting Maximum Link Utilization, Packet Loss, and Delay as metrics for comparison with OSPF.  

The main contributions of this work are the following:
\begin{enumerate}
\item The application and adaptation of Multi Agent Reinforcement Learning with Graph Neural Network for interconnection networks. 
\item A diagnosis of the conditions that cause MARL to fail and degrade routing performance.
\item A reproducible real-data evaluation pipeline of the multi-agent reinforcement learning algorithm using ns-3 simulator.
\end{enumerate}

\begin{figure}[htbp]
  \centering
  \includegraphics[width=\textwidth]{images/MARL Diagram.png}
  \caption{Multi-Agent Reinforcement Learning with GNN Illustration \cite{icon_neuralnetwork,icon_computer,icon_robot}}
  \label{fig:marl-diagram}
\end{figure}
